\section{A constant-distortion code for one prescribed length}\label{sec:code}

The short code in Lemma~\ref{lem:delimiter} is sufficient for both
tree lower bounds. Here we prove a stronger transfer result: a binary
block code with absolute distortion for every pair of words of one
prescribed length $M$, over an alphabet of size at most $M$. Its width
$M^2$ increases the word length polynomially. The proof combines
separation between segments of unequal codewords with an extraction of
original common subsequences from matches between equal codewords.

\begin{theorem}[Fixed-length binary conversion]\label{thm:804}
There is an absolute integer $M_0$ such that, for every integer
$M\geq M_0$ and every alphabet $\Gamma$ with $1\leq K=|\Gamma|\leq M$,
there exists a map $c:\Gamma\to\{0,1\}^{P}$ with $P=M^2$ whose
concatenation extension $\enc$ satisfies
\begin{equation}\label{eq:804}
 \frac P{804}\ED(x,y)
 \leq\ED(\enc(x),\enc(y))\leq P\ED(x,y)
 \qquad(x,y\in\Gamma^M).
\end{equation}
The encoded words have length $MP=M^3$. More precisely,
\begin{equation}\label{eq:402}
 r_{MP}(\enc(x),\enc(y))\geq\frac P{402}r_M(x,y)
 \qquad(x,y\in\Gamma^M).
\end{equation}
\end{theorem}

We first choose blocks whose long segments are separated whenever the
underlying letters differ. Separation of long intervals of distinct
binary codewords also appears in the inner code of Schulman and
Zuckerman \cite[Section~II-C, Condition~2 and Lemma~2]{SchulmanZuckerman1999}.
We prove the precise property needed here directly. The later argument
will handle matches between equal letters without a condition on
self-overlaps of a codeword.

\begin{lemma}[Separation of unequal-label segments]\label{lem:random-segments}
For sufficiently large integers $M$ and integers $1\leq K\leq M$,
there are $K$ binary words
of length $P=M^2$ such that any two length-$a$ segments of distinct
codewords satisfy $\LCS\leq0.99a$ whenever
\[
 a\geq a_0:=\lceil40\log_2(MP)\rceil
 =\lceil120\log_2M\rceil.
\]
\end{lemma}
\begin{proof}
Choose all $KP$ bits independently and uniformly. For two fixed
length-$a$ segments in distinct codewords, put $s=\lceil0.99a\rceil$.
Every prescribed increasing matching of size $s$ has probability
$2^{-s}$: its pairs use distinct independent bits on each side.
A union bound over the two position lists gives
\[
 \Prb[\LCS\geq s]\leq\binom as^2\,2^{-s}\leq2^{-a/2}.
\]
For the last estimate, let $u=a-s\leq0.01a$. If $u=0$, it is immediate.
Otherwise $\binom au\leq(ea/u)^u$, and
$u\log_2(ea/u)$ is increasing for $0<u<a$. Therefore
\[
 2\log_2\binom au-s
 \leq0.02a\log_2(100e)-0.99a<-0.81a<-a/2.
\]

There are at most $K^2P^3\leq M^2P^3$ choices of two codewords, a
segment length, and two starting positions. The probability that any
comparison of length at least $a_0$ fails is at most
\[
 M^2P^3\,2^{-a_0/2}
 \leq M^2P^3(MP)^{-20}<1.
\]
For large $M$, $a_0\leq P$. A realization avoiding every forbidden
comparison exists. Excluding $\LCS\geq\lceil0.99a\rceil$ implies
the stated weak inequality, including when $0.99a$ is an integer.
\end{proof}

\begin{proof}[Proof of Theorem~\ref{thm:804}]
Use codewords supplied by Lemma~\ref{lem:random-segments}. Fix any
increasing matching of $N_{\rm mat}$ bits between $\enc(x)$ and
$\enc(y)$, and put $\delta=MP-N_{\rm mat}$. Each encoded word has
exactly $\delta$ unmatched positions.

We first count matches between unequal original letters. Each matched
bit pair belongs to a pair $(i,j)$ of original block positions.
Along the matching, both block indices are nondecreasing. Every change
of pair increases $i+j$, so at most $2M-1$ block pairs occur. The
matches of one fixed pair are consecutive in this list: any pair
between two occurrences of $(i,j)$ must have both indices equal to
$i,j$.

Suppose a block pair with unequal letters contains $a$ matches.
Take the minimal contiguous spans containing those matched positions
in the two blocks. Let their lengths be $b,c$, and put
$g=(b-a)+(c-a)$. The spans for different block pairs are disjoint on
each side, even when the pairs share one block index. Moreover, no bit
matched under another pair can lie inside a span, by the consecutiveness
just proved. Thus $g$ counts globally unmatched positions, and the sum
of $g$ over all unequal-label pairs is at most $2\delta$.

If $a\geq a_0$, trim each span to a contiguous segment of length $a$.
This removes $g$ positions in total and destroys at most $g$ matched
pairs. The two segments retain at least $a-g$ matches, so
Lemma~\ref{lem:random-segments} gives $a-g\leq0.99a$. Hence
$g\geq0.01a$. Summing over these large pairs bounds their total matched
count by $100\sum g\leq200\delta$. The pairs with $a<a_0$ contribute
less than $2Ma_0$ matches. Consequently the number of unequal-label
matches is at most
\begin{equation}\label{eq:unequal-code}
 200\delta+2Ma_0.
\end{equation}

For the equal-label matches, form a bipartite multigraph whose two
vertex classes are the $M$ source and $M$ target block positions.
Each matched bit pair with equal original letters gives one edge.
Every vertex has degree at most $P$. The edges can be partitioned into
$P$ matchings. This is the classical bipartite edge-coloring theorem
of K\"{o}nig \cite{Konig1916}; we derive it here from
Hall's theorem. Add edges between deficient vertices, allowing parallel
edges, until the graph is $P$-regular. The total degree deficiency is
the same on the two sides, so this is possible. For any set $S$ of
left vertices, let $N(S)$ be its neighbor set. Counting incident edges
gives $P|S|\leq P|N(S)|$. Hall's theorem
\cite[Theorem~1]{Hall1935} supplies a perfect matching. Remove it and
repeat on the remaining regular graph, then discard all added edges.
This gives the required $P$ matchings of original edges.

Within one such matching, all block indices are distinct on both
sides. Since its edges came from the increasing bit matching, their
block pairs are strictly increasing in both indices. Their original
letters agree, so they form a common subsequence of $x,y$. Each of the
$P$ matchings therefore has at most $\LCS(x,y)$ edges, and the total
number of equal-label matches is at most $P\LCS(x,y)$.

Combine the two counts with $N_{\rm mat}=MP-\delta$:
\begin{equation}\label{eq:code-additive}
 P r_M(x,y)\leq201\delta+2Ma_0.
\end{equation}
Choose $M_0$ large enough that $2Ma_0\leq P/2$ for every $M\geq M_0$.
If $x\ne y$, then $r_M(x,y)\geq1$, and
\[
 201\delta\geq P r_M(x,y)-P/2
 \geq(P/2)r_M(x,y).
\]
Apply this to a longest encoded matching to obtain \eqref{eq:402}.
For $x=y$, that inequality is immediate. By \eqref{eq:deficit},
\[
 \ED(\enc(x),\enc(y))
 \geq r_{MP}(\enc(x),\enc(y))
 \geq\frac P{402}r_M(x,y)
 \geq\frac P{804}\ED(x,y).
\]
For the upper bound, extend $c$ to all finite words by concatenation.
Simulate each insertion or deletion in an optimal original edit script
by $P$ bit operations and each substitution by at most $P$ bit
substitutions. This proves \eqref{eq:804}; the intermediate word lengths
may be arbitrary.
\end{proof}

The segment property is a property of the chosen codebook, so the
theorem holds simultaneously for every pair in $\Gamma^M$. The choice
is tied to $M$: its guarantee is for one prescribed input length and
does not assert a logarithmic block width.

For the finite tree, $K\leq2m^k k^{5k}\leq M$ for all large $k$.
This theorem therefore also gives binary witnesses of length $M^3$
with distortion at least $(k/4)^k/(804C_1)$. The scale $P$ cancels
when an embedding is composed with the code. This alternative retains
an absolute distortion loss, at the cost of longer words than the
delimiter conversion used in Section~\ref{sec:finite}.
